\chapter{Results}
\label{chap:results}

This chapter reports what the agent actually did, question by question.
All figures come from two Easy-tier runs of 50 theorems each --- one against
NVIDIA Nemotron-3-Ultra-550B (run \texttt{20260909\_002719}), one against Groq
\texttt{gpt-oss-20b} (run \texttt{20260913\_165146}) --- with a step budget of 15.
Where a figure is drawn from the wider set of 20 run directories rather than from
these two runs alone, the text says so.
Medium and Hard tiers are unrun; \S\ref{sec:results_limits} states what that costs.

\section{RQ1 --- Capability ceiling}
\label{sec:results_rq1}

Table~\ref{tab:pass_rates} gives the two pass-rate variants defined in
Chapter~\ref{chap:problem_statement}.

\begin{table}[htbp]
\centering
\caption{Easy-tier pass rates, $n=50$ per run. The reachable denominator counts
only theorems the model was actually asked to prove and either proved or failed
to prove within budget.}
\label{tab:pass_rates}
\begin{tabular}{@{}lccc@{}}
\hline
\textbf{Model} & \textbf{Headline} & \textbf{Reachable-only} & \textbf{Reachable $n$} \\
\hline
Nemotron-3-Ultra-550B      & 7/50 = 14.0\% & 7/22 = 31.8\% & 22 \\
Groq \texttt{gpt-oss-20b}  & 8/50 = 16.0\% & 8/28 = 28.6\% & 28 \\
\hline
\end{tabular}
\end{table}

The gap between the two columns is the whole point.
On the headline figure the smaller model appears to lead by two percentage points;
on the reachable figure the larger model leads by three.
Neither difference is meaningful --- the Wilson intervals at these sample sizes
overlap almost completely --- but the fact that the ordering \emph{reverses}
under a defensible change of denominator is the clearest possible evidence that
the headline number is not measuring model capability.

What the headline denominator is measuring instead is visible in the result codes.
Of Nemotron's 50 attempts, 24 ended in \texttt{FAILED\_LLM\_ERROR}: the provider
returned errors and the model was never reached.
Of Groq's 50, 15 ended in \texttt{FAILED\_CONNECTIVITY\_OUTAGE} and 3 more in
\texttt{FAILED\_LLM\_ERROR}, the residue of free-tier rate limiting.
Between them the two runs lost 42 of 100 theorems to infrastructure.

\subsection{Elaboration failures}

Four theorems failed to elaborate in Lean before any tactic was proposed:
\texttt{exercise\_4\_4\_6a}, \texttt{exercise\_4\_5\_25}, \texttt{exercise\_6\_7},
and \texttt{exercise\_11\_4\_6b}.
The same four failed in both runs, which is the useful part of the observation:
an elaboration failure is a property of the dataset, not of the model, and it
reproduces exactly across providers.
These are statements whose Lean~4 rendering does not typecheck against current
Mathlib --- stale names, missing instances, or signatures that have drifted since
the problem set was transcribed.
They are excluded from the reachable denominator because no policy, however good,
could have proved them.

\subsection{Theorems proved}

Across both runs the agent proved nine distinct theorems:

\begin{center}
\begin{tabular}{@{}llc@{}}
\hline
\textbf{Theorem} & \textbf{Domain} & \textbf{Proved by} \\
\hline
\texttt{ALG\_E\_01} & Algebra          & both \\
\texttt{ALG\_E\_02} & Algebra          & both \\
\texttt{ALG\_E\_03} & Algebra          & both \\
\texttt{SET\_E\_01} & Set theory       & both \\
\texttt{SET\_E\_02} & Set theory       & both \\
\texttt{SET\_E\_03} & Set theory       & Nemotron only \\
\texttt{TOP\_E\_02} & Topology         & Groq only \\
\texttt{LIN\_E\_01} & Linear algebra   & both \\
\texttt{ABA\_E\_01} & Abstract algebra & Groq only \\
\hline
\end{tabular}
\end{center}

\texttt{LIN\_E\_01} and \texttt{ABA\_E\_01} appear in the raw logs under their
source names \texttt{exercise\_1\_3} and \texttt{exercise\_7\_1\_11}; they were
renamed on promotion into the problem set, and the union of the two runs' proved
sets is nine theorems rather than fifteen because five were proved twice.

The overlap is itself a result.
Seven of the nine were proved by both models, and the two singletons are not
obviously harder than the shared seven.
Two models three orders of magnitude apart in parameter count succeed on very
nearly the same set --- which suggests that on this tier the binding constraint
is not model scale but whether the theorem happens to fall to a single
well-known Mathlib lemma.

\section{RQ2 --- Failure mode taxonomy}
\label{sec:results_rq2}

Classification follows Chapter~\ref{chap:problem_statement}: five modes,
first-match-wins, with provider-caused failures excluded before classification.
That exclusion matters enormously here and is worth quantifying.
Classified without it, the 85 failed proofs across both runs yield 28
\textsc{Error Loop} cases --- but a run of identical \texttt{APIConnectionError}
messages satisfies the \textsc{Error Loop} test exactly, so most of those 28 are
outages wearing the costume of a reasoning failure.
Restricting to the 35 failures the model is actually responsible for
(\texttt{FAILED\_BUDGET\_EXHAUSTED}) gives:

\begin{center}
\begin{tabular}{@{}lrr@{}}
\hline
\textbf{Mode} & \textbf{Count} & \textbf{Share} \\
\hline
\textsc{No-Progress Loop}         &  0 &  0.0\% \\
\textsc{Error Loop}               &  6 & 17.1\% \\
\textsc{Hallucinated Lemma}       &  2 &  5.7\% \\
\textsc{Syntax Error Loop}        &  2 &  5.7\% \\
\textsc{Budget Exhaustion (Mixed)}& 25 & 71.4\% \\
\hline
\textbf{Total}                    & 35 & 100.0\% \\
\hline
\end{tabular}
\end{center}

The headline is the zero, and \S\ref{sec:results_rq4} takes it up properly.

The rest of the distribution is a tail behind one dominant bucket.
\textsc{Budget Exhaustion (Mixed)} is not a pathology --- it is the absence of one.
It means the agent tried fifteen different things, none of them repeated often
enough or wrongly enough to trip a named mode, and then ran out of steps.
That nearly three quarters of genuine failures land there is a negative result
about the taxonomy as much as about the agent: the five modes were designed
expecting the agent to get \emph{stuck} in a characteristic way, and mostly it
does not get stuck, it simply does not know the answer.

\textsc{Hallucinated Lemma} deserves a caveat in the other direction.
Two confirmed cases is a floor, not an estimate.
The detector fires only when Lean says ``unknown identifier'' \emph{and} the
tactic names a capitalised non-stdlib symbol; a model that invents a lowercase
lemma name, or whose invented name happens to exist with the wrong signature,
produces a type error rather than an unknown identifier and lands in
\textsc{Budget Exhaustion} instead. The true rate is higher than 5.7\%.

\section{RQ3 --- Step economy}
\label{sec:results_rq3}

Two different quantities can both be called ``steps to proof'', and they give
very different answers. Across all 38 proofs in all 20 run directories:

\begin{center}
\begin{tabular}{@{}lccc@{}}
\hline
\textbf{Measure} & \textbf{Min} & \textbf{Median} & \textbf{Max} \\
\hline
Attempts made (all tactics proposed)     & 1 & 1 & 9 \\
Proof length (tactics that survived)     & 1 & 1 & 3 \\
\hline
\end{tabular}
\end{center}

\emph{Attempts made} counts every tactic the agent proposed, failures included:
20 of 38 proofs closed on the first attempt, 33 of 38 within four, and the
longest took nine.
\emph{Proof length} counts only the steps whose outcome was \texttt{advanced} or
\texttt{proved} --- the tactics that actually appear in the compiled proof, since
a tactic that errored left the proof state untouched.
By that measure no proof exceeds three tactics and 34 of 38 are a single tactic.

The distinction is not pedantry: it is the difference between what the agent did
and what it produced. The proofs this agent finds are one-liners that it
sometimes takes several guesses to hit.

For early stopping the operative number is the first one, and it is encouraging.
No proof in any run closed after step~9, and the marginal return past step~4 is
five proofs out of 38. With a budget of 15, roughly two thirds of the compute
spent on eventually-successful proofs is spent after the point where success has
become unlikely --- and all of the compute on the 25 budget-exhausted failures is
spent there. Cutting the budget to 6 would have cost two proofs and saved a
large majority of the wasted steps.

\section{RQ4 --- History utilisation}
\label{sec:results_rq4}

\subsection{The hypothesis was wrong}

The study was designed around the expectation that \textsc{No-Progress Loop}
would be the dominant failure mode, on the order of 40\%: the agent would propose
a tactic, Lean would accept it without changing the goal, the agent would fail to
notice and propose it again, and the episode would spin in place until the budget
ran out.

It occurred zero times. Not rarely --- never, in either run, under either model,
and not in any of the 20 run directories.

The mechanism is straightforward once observed. The failure mode presupposes a
silent failure: a tactic that Lean accepts while leaving the goal untouched.
In practice the agent's bad tactics do not fail silently, they fail loudly.
Lean returns an error message, that message is appended to the history, and the
next prompt contains it. The agent is never in the epistemic position the
hypothesis assumed. Whatever else it does wrong, it is never left guessing
whether its last move did anything.

This is a genuine negative result and it redirects the improvement effort.
Work aimed at making the agent notice stagnation --- goal-state diffing,
explicit no-progress detection, forced tactic diversity --- addresses a condition
that does not arise. The agent's problem is not that it fails to notice failure.
It is that noticing failure does not tell it what to do next.

\subsection{Repetition}

The agent repeats itself, but not blindly.
Pooled over all steps, 36.4\% of Nemotron's steps and 32.6\% of Groq's proposed a
tactic that had already appeared earlier in the same proof (34.9\% combined,
29.6\% across all 20 run directories).
The distribution is skewed: the median proof repeats little, while a handful
repeat heavily --- two Groq proofs repeated on 11 of 15 steps.

Repetition after an error is not the same defect as repetition after no-progress.
The agent has been told the tactic failed and why, and tries it anyway --- often
with a small perturbation that does not address the error.
This is a failure to use the feedback, not a failure to receive it.

\subsection{Within-episode learning}

If the agent learned from its history, the \texttt{advanced} rate would rise with
step number. It falls:

\begin{center}
\begin{tabular}{@{}lcccccc@{}}
\hline
\textbf{Step} & 1 & 2 & 3 & 4 & 5 & 6--8 \\
\hline
\textbf{Advanced rate} & 23.9\% & 11.9\% & 4.9\% & 16.7\% & 6.2\% & 0.0\% \\
\hline
\end{tabular}
\end{center}

\noindent (Nemotron run; the denominator shrinks as proofs terminate.)
The first tactic is by far the most likely to work, and the curve decays to zero
by step~6. The agent front-loads its best guess and then degrades.

Goal-state size tells the same story. Measured in characters of the goal shown to
the model, proofs that eventually succeeded went 64, 66, 71, 71 at steps 1, 3, 6
and 10; proofs that failed went 86, 87, 87, 87.
Both rows are flat or rising. The agent is not decomposing the problem into
smaller pieces --- it is restating it. A successful proof is one where the first
or second guess happened to close a goal that was small to begin with.

\section{RQ5 --- Tactic distribution}
\label{sec:results_rq5}

The single most striking result in this chapter is which tactic each model used
most often.

For Nemotron, over 231 tactic uses spanning 124 distinct tactics, the most
frequent tactic was \texttt{sorry} --- 19 uses, 8.2\% of everything it proposed.
\texttt{sorry} is Lean's placeholder for an omitted proof. It closes the goal
by assumption and proves nothing; the harness scores it as a failure, and its
precision over those 19 uses is 0.00.

For Groq the same behaviour appears in a different costume and at greater
frequency: \texttt{done} (20 uses) and \texttt{admit} (20 uses) top the table,
followed by \texttt{exact sorry} (7) and \texttt{exact (by admit)} (3).
Together those four are 50 of 141 uses --- 35.5\% of everything the smaller model
proposed was a way of giving up.

The natural reading is that the model loses track of what it is being asked to do.
The prompt contains a theorem statement ending in \texttt{by sorry} --- the
benchmark's own placeholder --- and the model, pattern-matching on a
familiar shape, completes the file as a source-code exercise rather than
proposing the next tactic in a live proof state. It is not conceding;
it is answering a different question.

This is an important target because it is cheap to fix and it is not a capability
limit. Roughly a third of the smaller model's step budget is being spent on
tactics that cannot possibly succeed. A prompt-level change, or a harness-level
rejection of surrender tactics with a re-ask, would recover that budget without
any improvement in mathematical ability.

Among tactics that did work, the pattern is narrow. Groq's \texttt{simp} closed
3 of 5 uses (precision 0.60), the only tactic in either run to clear the
three-use minimum with a positive rate. The proofs that succeeded were closed by
single applications of exactly the right Mathlib lemma --- \texttt{rw [neg\_neg]},
\texttt{exact sq\_eq\_one\_iff.mp hx}, \texttt{exact h1.trans h2} --- which is
consistent with RQ3's finding that successful proofs are one-liners.

\section{Limitations}
\label{sec:results_limits}

\paragraph{Sample size.}
Fifty theorems per run, of which 22 and 28 were reachable. Every rate in this
chapter carries a Wilson interval wide enough to accommodate most competing
hypotheses. No claim of a difference between the two models is supported.

\paragraph{Medium and Hard tiers are unrun.}
The tier labels are inherited from the source problem set and describe human
difficulty; whether they predict agent difficulty is exactly RQ1's second half
and it remains unanswered. The runs are gated on provider reliability rather than
on compute: an Easy-50 run that loses 42\% of its theorems to outages would lose
proportionally more of a 378-theorem Medium tier, and the resulting headline
figure would be uninterpretable. Reporting Easy-only is the honest consequence of
that, not an omission.

\paragraph{The failure taxonomy is coarse.}
71\% of genuine failures land in the residual bucket. The taxonomy successfully
falsified its central hypothesis (\S\ref{sec:results_rq4}) but it does not
discriminate well among the failures that remain, and
\textsc{Hallucinated Lemma} is a known undercount.

\paragraph{Single-sample, zero-shot.}
One attempt per theorem, one tactic per step, no temperature sweep, no
self-consistency. These results bound what a minimal agent does, not what these
models can do.
